% Figure 4.2 --- three residual gaps of TOPORET, closed by DOTS
\begin{figure}[!tb]
\centering
\resizebox{\linewidth}{!}{%
\begin{tikzpicture}[
  font=\footnotesize,
  cell/.style={draw, rounded corners=2.5pt, align=center, inner sep=4pt,
               minimum height=1.55cm, text width=4.0cm},
  arr/.style={-{Latex[length=2mm]}, thick}
]
\def\dx{4.9}\def\dy{2.0}
\node[font=\small\bfseries, text=failred] at (-\dx,0) {Gap left by TOPORET};
\node[font=\small\bfseries, text=navy] at (0,0) {What DOTS adds};
\node[font=\small\bfseries, text=okgreen] at (\dx,0) {Measured effect};

\node[cell, fill=failf, draw=failred] (a1) at (-\dx,-1.2) {\textbf{A -- domain mismatch}\\ImageNet backbone sees\\textures, not the retina};
\node[cell, fill=black!4, draw=black!50] (a2) at (0,-1.2) {RETFound ViT-L/16\\pre-trained on $1.6$M\\retinal images};
\node[cell, fill=okf, draw=okgreen] (a3) at (\dx,-1.2) {\textbf{$+4.1$ pp QWK}\\Gap A closed};

\node[cell, fill=failf, draw=failred] (b1) at (-\dx,-1.2-\dy) {\textbf{B -- ordinal violations}\\CORN: $2.1\%$ of rankings\\are logically impossible};
\node[cell, fill=black!4, draw=black!50] (b2) at (0,-1.2-\dy) {CORAL monotone bias\\(algebraic constraint,\\not a loss penalty)};
\node[cell, fill=okf, draw=okgreen] (b3) at (\dx,-1.2-\dy) {\textbf{$0\%$ violations}\\Gap B closed};

\node[cell, fill=failf, draw=failred] (c1) at (-\dx,-1.2-2*\dy) {\textbf{C -- safety unmet}\\no right to abstain;\\Sev-NR $=3.4\%$};
\node[cell, fill=black!4, draw=black!50] (c2) at (0,-1.2-2*\dy) {dynamic GAT $+$\\entropy abstention\\($\tau_U=0.73$ nats)};
\node[cell, fill=okf, draw=okgreen] (c3) at (\dx,-1.2-2*\dy) {\textbf{Sev-NR $=1.9\%$}\\target met at the point\\estimate (CI up to $2.4\%$)};

\foreach \a/\b in {a1/a2,a2/a3,b1/b2,b2/b3,c1/c2,c2/c3} \draw[arr] (\a) -- (\b);
\end{tikzpicture}}
\caption[The three residual gaps of TOPORET]{The three residual gaps of TOPORET, measured on its ordinal variant under the protocol of this chapter. Read
each row left to right: the problem, the DOTS fix, the measured effect.
\Cref{sec:gap-a,sec:gap-b,sec:gap-c} take the rows in order.}
\label{fig:gaps}
\end{figure}
